% =============================================================
% บทที่ 2 คู่มือการใช้งาน
% เนื้อหาตรวจจากเว็บจริง https://dinokhonkaen.app และซอร์สโค้ดใน main repo/frontend
% ภาพหน้าจอ: images/screens/<ชื่อ>.jpg|png  (ถ้ายังไม่มีไฟล์ จะแสดงกรอบแดงแทน)
% =============================================================

\section{คู่มือการใช้งาน}

``ไดโน'' เว็บแอปพลิเคชันศูนย์กลางของการท่องเที่ยวจังหวัดขอนแก่น เข้าใช้งานได้ที่ https://dinokhonkaen.app แบ่งผู้ใช้งานออกเป็น 3 กลุ่ม ได้แก่

\begin{numlist}
\item ผู้ใช้งานทั่วไป (Guest)
\item สมาชิก (Tourist)
\item ผู้ดูแลระบบ (Admin)
\end{numlist}

% -------------------------------------------------------------
\subsection{เข้าใช้งานกรณีเป็นผู้ใช้งานทั่วไป}
% -------------------------------------------------------------

แถบเมนูผู้ใช้งานทั่วไป อยู่ด้านบนของทุกหน้า ดังภาพที่ \ref{fig:guest-home}

\screenshot{guest-home}{หน้าแรกและแถบเมนูผู้ใช้งานทั่วไป}

จากภาพที่ \ref{fig:guest-home} รายละเอียดแถบเมนูมีดังนี้

\begin{numlist}
\item หน้าแรก แสดงช่องค้นหา กิจกรรมและเทศกาลที่กำลังจะมาถึง และสถานที่แนะนำ
\item สถานที่ท่องเที่ยว แสดงรายการสถานที่ทั้งหมดพร้อมตัวกรองตามหมวดหมู่
\item กิจกรรม \& เทศกาล แสดงรายการกิจกรรมและเทศกาลในขอนแก่น
\item วางแผนทริป AI ใช้สร้างแผนการเดินทางให้อัตโนมัติ
\item พอยท์สะสม แสดงสถานที่ร่วมรายการและของรางวัล
\item เข้าสู่ระบบ และสมัครสมาชิก อยู่มุมขวาบน
\end{numlist}

เมื่อเปิดเว็บไซต์ครั้งแรก ระบบอาจแสดงหน้าต่างต้อนรับจากน้องไดโน สามารถกดปุ่ม ``เริ่มสำรวจขอนแก่น'' หรือ ``ไม่แสดงอีก 7 วัน'' เพื่อปิดได้ และมีปุ่มรูปน้องไดโนมุมขวาล่างของทุกหน้าสำหรับเปิดแชตบอต

\subsubsection{เมนูสมัครสมาชิก กรณีเป็นผู้ใช้งานทั่วไป}

การสมัครสมาชิกทำได้โดยคลิกปุ่ม ``สมัครสมาชิก'' มุมขวาบนของหน้าเว็บ ระบบจะแสดงแบบฟอร์มสมัครสมาชิกดังภาพที่ \ref{fig:guest-signup-form}

\screenshot{guest-signup-form}{หน้าฟอร์มสมัครสมาชิก}

จากภาพที่ \ref{fig:guest-signup-form} ผู้ใช้กรอกข้อมูล ได้แก่ คำนำหน้า ชื่อ นามสกุล เพศ วันเกิด อาชีพ อีเมล เบอร์โทรศัพท์ รหัสผ่านและยืนยันรหัสผ่าน และที่อยู่ปัจจุบัน (จังหวัด อำเภอ ตำบล) โดยรหัสผ่านต้องมีความยาวอย่างน้อย 8 ตัวอักษร ประกอบด้วยตัวพิมพ์ใหญ่ ตัวพิมพ์เล็ก และตัวเลข ระบบแสดงผลการตรวจสอบเงื่อนไขทีละข้อใต้ช่องรหัสผ่าน จากนั้นต้องทำเครื่องหมายยินยอมให้เก็บและใช้ข้อมูลส่วนบุคคลตามนโยบายความเป็นส่วนตัว แล้วกดปุ่ม ``สมัครสมาชิก''

หลังสมัครสำเร็จ ระบบจะส่งอีเมลยืนยันไปยังที่อยู่อีเมลที่กรอก ผู้ใช้ต้องคลิกลิงก์ในอีเมลเพื่อยืนยันบัญชี \fillin{ยืนยันกับระบบจริงว่าต้องยืนยันอีเมลก่อนเข้าสู่ระบบหรือไม่}

\subsubsection{เมนูเข้าสู่ระบบ กรณีเป็นผู้ใช้งานทั่วไป}

ผู้ใช้ที่มีบัญชีแล้วคลิกเมนู ``เข้าสู่ระบบ'' มุมขวาบน ระบบจะแสดงแบบฟอร์มเข้าสู่ระบบดังภาพที่ \ref{fig:guest-login-form}

\screenshot{guest-login-form}{หน้าฟอร์มเข้าสู่ระบบ}

ผู้ใช้กรอกอีเมลและรหัสผ่านที่ลงทะเบียนไว้ แล้วกดปุ่ม ``เข้าสู่ระบบ'' หากยังไม่มีบัญชีสามารถคลิก ``สมัครสมาชิก'' ด้านล่างปุ่ม เพื่อไปยังหน้าสมัครสมาชิกได้ทันที

หากลืมรหัสผ่าน ให้คลิก ``ลืมรหัสผ่าน?'' เพื่อเข้าสู่หน้าดังภาพที่ \ref{fig:forgot-password} กรอกอีเมลที่ใช้สมัครแล้วกด ``ส่งลิงก์รีเซ็ตรหัสผ่าน'' ระบบจะส่งลิงก์สำหรับตั้งรหัสผ่านใหม่ไปยังอีเมลนั้น

\screenshot{forgot-password}{หน้าลืมรหัสผ่าน}

\subsubsection{เมนูสถานที่ท่องเที่ยว กรณีเป็นผู้ใช้งานทั่วไป}

ผู้ใช้เรียกดูสถานที่ได้จากเมนู ``สถานที่ท่องเที่ยว'' ระบบจะแสดงรายการสถานที่ในรูปแบบการ์ด พร้อมช่องค้นหาและตัวกรองหมวดหมู่ ได้แก่ ทั้งหมด คาเฟ่ วัด พิพิธภัณฑ์ สวนสาธารณะ อุทยานแห่งชาติ ตลาด สถานที่ท่องเที่ยว และร้านอาหาร ดังภาพที่ \ref{fig:places-list}

\screenshot{places-list}{หน้ารายการสถานที่ท่องเที่ยว}

เมื่อคลิกการ์ดสถานที่ ระบบจะแสดงหน้ารายละเอียดดังภาพที่ \ref{fig:place-detail} ซึ่งประกอบด้วยหมวดหมู่ คะแนนและช่วงราคา แกลเลอรีรูปภาพ สรุปโดย Dino AI สิ่งอำนวยความสะดวก และรีวิวจากผู้เยี่ยมชม พร้อมปุ่ม ``เปิดเส้นทาง Google Maps'' และปุ่ม ``โทรติดต่อ''

\screenshot{place-detail}{หน้ารายละเอียดสถานที่ท่องเที่ยว}

\subsubsection{เมนูกิจกรรมและเทศกาล กรณีเป็นผู้ใช้งานทั่วไป}

เมนู ``กิจกรรม \& เทศกาล'' แสดงรายการกิจกรรมและเทศกาลในขอนแก่น พร้อมช่องค้นหาและตัวกรองสถานะ ได้แก่ ทั้งหมด กำลังจัดอยู่ เร็ว ๆ นี้ จบแล้ว และยกเลิก ดังภาพที่ \ref{fig:events-list}

\screenshot{events-list}{หน้ารายการกิจกรรมและเทศกาล}

เมื่อคลิกกิจกรรม ระบบจะแสดงรายละเอียด ได้แก่ รูปภาพ คำอธิบาย วันจัดงาน สถานที่จัดงาน ค่าเข้าชม ผู้จัดงาน กลุ่มที่เหมาะสม และแผนที่ ดังภาพที่ \ref{fig:event-detail} หากพบข้อมูลไม่ถูกต้อง สามารถกดปุ่ม ``แจ้งข้อมูลไม่ถูกต้อง'' ใต้รูปภาพเพื่อส่งให้ผู้ดูแลระบบตรวจสอบได้

\screenshot{event-detail}{หน้ารายละเอียดกิจกรรม}

\subsubsection{เมนูแชตบอต กรณีเป็นผู้ใช้งานทั่วไป}

ผู้ใช้คลิกปุ่มรูปน้องไดโนมุมขวาล่างเพื่อเปิดหน้าต่างแชตบอต พิมพ์คำถามเกี่ยวกับสถานที่ อาหาร หรือเทศกาลในช่องด้านล่าง แล้วกดปุ่มส่ง ดังภาพที่ \ref{fig:chatbot}

\screenshot[0.9]{chatbot}{หน้าต่างแชตบอตน้องไดโน}

ระบบจะค้นหาข้อมูลที่เกี่ยวข้องจากฐานความรู้และตอบกลับแบบทยอยปรากฏทีละข้อความ หากไม่พบข้อมูลที่เกี่ยวข้อง แชตบอตจะแจ้งว่าไม่มีข้อมูลแทนการตอบโดยไม่มีข้อมูลรองรับ ในกรณีที่ปรากฏข้อความ ``ระบบแชทขัดข้องชั่วคราว'' ให้ลองใหม่ภายหลัง

\subsubsection{เมนูวางแผนทริป AI กรณีเป็นผู้ใช้งานทั่วไป}

ผู้ใช้เลือกเมนู ``วางแผนทริป AI'' เพื่อสร้างแผนการเดินทางผ่านแบบฟอร์ม 4 ขั้นตอน โดยมีแถบแสดงความคืบหน้าด้านบน ขั้นตอนที่ 1 วันเดินทาง ดังภาพที่ \ref{fig:trip-step1}

\screenshot{trip-step1}{หน้าแบบฟอร์มสร้างแผนการเดินทาง ขั้นตอนที่ 1 วันเดินทาง}

จากภาพที่ \ref{fig:trip-step1} ผู้ใช้เลือกช่วงเวลาเดินทางจากปุ่มลัด ได้แก่ วันนี้ (เดย์ทริป) 2 วัน 3 วัน หรือสุดสัปดาห์นี้ หรือกำหนดวันเริ่มต้นและวันสิ้นสุดเอง ระบุที่พัก โดยพิมพ์ค้นหาหรือกด ``ตำแหน่งปัจจุบัน'' และระบุสถานที่ที่ต้องการไปแน่นอน (ถ้ามี) หากไม่ระบุที่พัก ระบบจะใช้ใจกลางเมืองขอนแก่นแทน จากนั้นกด ``ถัดไป''

ขั้นตอนที่ 2 ความสนใจ ผู้ใช้เลือกได้หลายข้อ ได้แก่ ธรรมชาติ วัฒนธรรม/ศาสนา ไดโนเสาร์ อาหารพื้นถิ่น คาเฟ่ ช้อปปิ้ง/หัตถกรรม และครอบครัว ดังภาพที่ \ref{fig:trip-step2}

\screenshot{trip-step2}{หน้าแบบฟอร์มสร้างแผนการเดินทาง ขั้นตอนที่ 2 ความสนใจ}

ขั้นตอนที่ 3 งบประมาณ ผู้ใช้เลือกงบประมาณโดยรวม (ประหยัด ปานกลาง หรูหรา) จังหวะการเที่ยว (สายชิลล์ กำลังดี หรือสายลุย) และขอบเขตพื้นที่ (เมือง หรือทั่วขอนแก่นซึ่งรวมสถานที่รอบนอก เช่น ภูเวียง อุบลรัตน์) ดังภาพที่ \ref{fig:trip-step3}

\screenshot{trip-step3}{หน้าแบบฟอร์มสร้างแผนการเดินทาง ขั้นตอนที่ 3 งบประมาณ}

ขั้นตอนที่ 4 ช่วงเวลา ผู้ใช้กำหนดเวลาเริ่มและเวลาสิ้นสุดของแต่ละวัน (ค่าเริ่มต้น 09:00 ถึง 18:00) แล้วกดปุ่ม ``สร้างแผนการเดินทาง'' ดังภาพที่ \ref{fig:trip-step4}

\screenshot{trip-step4}{หน้าแบบฟอร์มสร้างแผนการเดินทาง ขั้นตอนที่ 4 ช่วงเวลา}

เมื่อกดสร้างแผน ระบบแสดงหน้าสถานะกำลังประมวลผล ดังภาพที่ \ref{fig:trip-loading} จากนั้นแสดงแผนการเดินทางรายวันพร้อมแผนที่เส้นทาง ดังภาพที่ \ref{fig:trip-result}

\screenshot[0.7]{trip-loading}{หน้าสถานะกำลังประมวลผลแผนการเดินทาง}

\screenshot{trip-result}{หน้าแสดงผลแผนการเดินทาง}

ในหน้าผลลัพธ์ ผู้ใช้สามารถพิมพ์สิ่งที่อยากปรับ แล้วกดปุ่ม ``ปรับแผนการเดินทางอีกครั้ง'' เพื่อให้ระบบสร้างแผนใหม่ ทั้งนี้ ผู้ใช้งานทั่วไปที่ไม่ได้เข้าสู่ระบบจะไม่สามารถบันทึกแผนได้ แผนจะหายเมื่อปิดหน้า จึงควรเข้าสู่ระบบก่อนสร้างแผนเพื่อบันทึกไว้ในเมนู ``ทริปของฉัน''

\subsubsection{เมนูพอยท์สะสม กรณีเป็นผู้ใช้งานทั่วไป}

เมนู ``พอยท์สะสม'' สำหรับผู้ที่ยังไม่เข้าสู่ระบบ แสดงบัตรสะสมพอยท์ที่ระบุสถานที่ร่วมรายการพร้อมจำนวนพอยท์ที่ได้รับจากแต่ละที่ ขั้นตอนการสะสม 3 ขั้นตอน และรายการของรางวัลที่แลกได้ ดังภาพที่ \ref{fig:points-guest}

\screenshot{points-guest}{หน้าพอยท์สะสม กรณีผู้ใช้งานทั่วไป}

ผู้ใช้สามารถกด ``สมัครสมาชิกฟรี'' หรือ ``เข้าสู่ระบบ'' เพื่อเริ่มสะสมพอยท์ได้

% -------------------------------------------------------------
\subsection{เข้าใช้งานกรณีเข้าสู่ระบบของสมาชิก}
% -------------------------------------------------------------

เมื่อเข้าสู่ระบบแล้ว แถบเมนูจะเพิ่มเมนู ``ทริปของฉัน'' และแทนที่ปุ่มเข้าสู่ระบบด้วยเมนูผู้ใช้มุมขวาบน ดังภาพที่ \ref{fig:member-navbar}

\screenshot{member-navbar}{แถบเมนูสมาชิก}

\subsubsection{เมนูพอยท์สะสม กรณีเป็นสมาชิก}

หน้า ``พอยท์สะสมของคุณ'' แสดงยอดพอยท์ที่มีอยู่ ปุ่มสแกน QR Code รายการสถานที่ที่มี QR สำหรับรับพอยท์ และรายการของรางวัล ดังภาพที่ \ref{fig:points-member}

\screenshot{points-member}{หน้าพอยท์สะสมของสมาชิก}

\subsubsection{การสแกน QR Code เพื่อสะสมพอยท์ กรณีเป็นสมาชิก}

เมื่อไปถึงสถานที่ร่วมรายการ สมาชิกสแกนป้าย QR Code ด้วยปุ่มสแกนในหน้าพอยท์สะสม หรือใช้กล้องของโทรศัพท์สแกนโดยตรงก็ได้ ระบบจะขออนุญาตใช้กล้อง ดังภาพที่ \ref{fig:qr-scan}

\screenshot[0.7]{qr-scan}{หน้าต่างสแกน QR Code}

เมื่อสแกนสำเร็จ ระบบเพิ่มพอยท์เข้าบัญชีทันทีและแสดงผลการสแกน QR Code แต่ละจุดรับพอยท์ได้ครั้งเดียวต่อบัญชี หากสแกนซ้ำหรือ QR Code ไม่ถูกต้อง ระบบจะแสดงข้อความ ``สแกนไม่สำเร็จ'' พร้อมเหตุผล

\subsubsection{การแลกของรางวัล กรณีเป็นสมาชิก}

การแลกของรางวัลดำเนินการที่เคาน์เตอร์ โดยสมาชิกแจ้งชื่อหรือเบอร์โทรศัพท์ที่ใช้สมัครกับเจ้าหน้าที่ เจ้าหน้าที่จะหักพอยท์และมอบของรางวัลให้ผ่านหน้าผู้ดูแลระบบตามหัวข้อที่ \ref{sec:admin-guide}

\subsubsection{เมนูทริปของฉัน กรณีเป็นสมาชิก}

เมนู ``ทริปของฉัน'' แสดงแผนทริปที่บันทึกไว้ทั้งหมด ดังภาพที่ \ref{fig:my-trips} ผู้ใช้สามารถค้นหาตามชื่อทริป เปลี่ยนชื่อ เพิ่มในรายการโปรด ทำสำเนา และลบแผนได้ รวมถึงกด ``+ สร้างแผนใหม่'' เพื่อสร้างแผนเพิ่ม

\screenshot{my-trips}{หน้าทริปของฉัน}

\subsubsection{เมนูโปรไฟล์ กรณีเป็นสมาชิก}

ผู้ใช้เปิดหน้าโปรไฟล์ได้จากเมนูผู้ใช้มุมขวาบน ประกอบด้วยแท็บ 5 แท็บ ได้แก่ ภาพรวม ข้อมูล ประวัติ กิจกรรมของฉัน และความปลอดภัย ดังภาพที่ \ref{fig:profile-overview}

\screenshot{profile-overview}{หน้าโปรไฟล์ แท็บภาพรวม}

\begin{numlist}
\item ภาพรวม แสดงข้อมูลสรุปของผู้ใช้ ยอดพอยท์ และปุ่มลัด ``สแกน QR สะสมพอยท์'' และ ``ดูประวัติพอยท์''
\item ข้อมูล ใช้แก้ไขข้อมูลส่วนตัว เช่น คำนำหน้า ชื่อ นามสกุล วันเกิด อาชีพ เบอร์โทรศัพท์ และที่อยู่ แล้วกด ``บันทึกการเปลี่ยนแปลง'' ดังภาพที่ \ref{fig:profile-info}
\item ประวัติ แสดงประวัติพอยท์ทั้งการสแกน QR การแลกรางวัล และการปรับพอยท์ พร้อมตัวกรอง ดังภาพที่ \ref{fig:profile-history}
\item กิจกรรมของฉัน ใช้ส่งกิจกรรมใหม่เข้าระบบด้วยปุ่ม ``+ ส่งกิจกรรมใหม่'' และติดตามสถานะรออนุมัติ อนุมัติแล้ว หรือไม่อนุมัติ ดังภาพที่ \ref{fig:profile-events}
\item ความปลอดภัย ใช้เปลี่ยนรหัสผ่าน เปลี่ยนอีเมล และลบบัญชี ดังภาพที่ \ref{fig:profile-security}
\end{numlist}

\screenshot{profile-info}{หน้าโปรไฟล์ แท็บข้อมูล}

\screenshot{profile-history}{หน้าโปรไฟล์ แท็บประวัติ}

\screenshot{profile-events}{หน้าโปรไฟล์ แท็บกิจกรรมของฉัน}

\screenshot{profile-security}{หน้าโปรไฟล์ แท็บความปลอดภัย}

% -------------------------------------------------------------
\subsection{เข้าใช้งานกรณีเข้าสู่ระบบของผู้ดูแลระบบ}
\label{sec:admin-guide}
% -------------------------------------------------------------

\subsubsection{เมนูเข้าสู่ระบบ กรณีเป็นผู้ดูแลระบบ}

ผู้ดูแลระบบเข้าสู่ระบบที่ https://dinokhonkaen.app/admin/login หรือคลิกลิงก์ ``สำหรับผู้ดูแลระบบ'' ที่ส่วนท้ายของเว็บไซต์ ดังภาพที่ \ref{fig:admin-login} จากนั้นกรอกอีเมลและรหัสผ่านของบัญชีผู้ดูแลระบบแล้วกด ``เข้าสู่ระบบ'' บัญชีผู้ดูแลระบบสร้างโดยนักพัฒนาตามหัวข้อที่ \ref{sec:create-admin}

\screenshot[0.7]{admin-login}{หน้าเข้าสู่ระบบสำหรับผู้ดูแลระบบ}

\subsubsection{แถบเมนูและแดชบอร์ด}

หลังเข้าสู่ระบบ ระบบแสดงแถบเมนูด้านข้างที่จัดเป็นกลุ่ม ดังนี้ ภาพรวม (แดชบอร์ด) เนื้อหา (สถานที่ กิจกรรม ฐานความรู้แชทบอท) คิวตรวจสอบ (รายงานข้อมูล คำขอกิจกรรม) พอยท์และรางวัล (QR Code ของรางวัล แลกที่เคาน์เตอร์) และผู้ใช้และการใช้งาน (ผู้ใช้ ทริป ความพึงพอใจ) ซึ่งชื่อเมนูตรงกับที่แสดงในระบบ ส่วนแดชบอร์ดแบ่งเป็น 3 แท็บ ได้แก่ ภาพรวม ผู้ใช้และรางวัล และแผนทริป โดยแท็บภาพรวมแสดงจำนวนสถานที่ กิจกรรม ผู้ใช้ทั้งหมด ผู้สมัครใหม่ และการสแกน QR Code ใน 7 วัน รวมถึงจำนวนฐานความรู้ QR Code ของรางวัล และพอยท์ที่แจกไปแล้ว นอกจากนี้ยังมีส่วนอัปเดตดัชนีค้นหาของแชตบอตและเวอร์ชันของระบบ ดังภาพที่ \ref{fig:admin-dashboard}

\screenshot{admin-dashboard}{หน้าแดชบอร์ดผู้ดูแลระบบ}

\subsubsection{เมนูสถานที่}

ผู้ดูแลระบบจัดการข้อมูลสถานที่ได้จากเมนู ``สถานที่'' ดังภาพที่ \ref{fig:admin-places} โดยเพิ่มสถานที่ใหม่ แก้ไขข้อมูล ลบสถานที่ และนำเข้า/ซิงก์ข้อมูลจาก Google Places ได้ ฟอร์มเพิ่มและแก้ไขข้อมูลสถานที่แสดงดังภาพที่ \ref{fig:admin-place-form} ในการซิงก์ ผู้ดูแลระบบสามารถล็อกฟิลด์ที่แก้เองไว้ได้ ค่าใหม่จาก Google ของฟิลด์ที่ล็อกจะถูกพักไว้ให้ตรวจสอบก่อน และการซิงก์แต่ละสถานที่มีค่าใช้จ่ายตามการเรียก Google Places API

\screenshot{admin-places}{หน้าจัดการสถานที่}

\screenshot{admin-place-form}{หน้าฟอร์มเพิ่มและแก้ไขข้อมูลสถานที่}

\subsubsection{เมนูกิจกรรม}

เมนู ``กิจกรรม'' ใช้เพิ่ม แก้ไข และลบกิจกรรมและเทศกาล ดังภาพที่ \ref{fig:admin-events} ในฟอร์มกิจกรรมมีเครื่องมือช่วยดึงข้อมูลจากข้อความโพสต์ด้วย AI ให้ผู้ดูแลระบบวางข้อความโพสต์ แล้วระบบจะกรอกฟอร์มให้อัตโนมัติ ดังภาพที่ \ref{fig:admin-event-extract} ผู้ดูแลระบบควรตรวจสอบข้อมูลที่ AI กรอกให้ก่อนกดบันทึกทุกครั้ง

\screenshot{admin-events}{หน้าจัดการกิจกรรม}

\screenshot{admin-event-extract}{หน้าฟอร์มกิจกรรมพร้อมเครื่องมือดึงข้อมูลจากข้อความโพสต์ด้วย AI}

\subsubsection{เมนูฐานความรู้แชตบอต}

ผู้ดูแลระบบเพิ่ม แก้ไข ปักหมุด และเปิด/ปิดการใช้งานเอกสารความรู้ที่แชตบอตใช้อ้างอิง ดังภาพที่ \ref{fig:admin-knowledge} หลังเพิ่มหรือแก้ไขข้อมูล ต้องให้ระบบสร้าง Embedding ใหม่ เพื่อให้แชตบอตค้นพบข้อมูลล่าสุด โดยกดปุ่ม ``อัปเดตดัชนีค้นหา'' ในการ์ด ``ดัชนีค้นหาแชตบอท (Embedding)'' บนหน้าแดชบอร์ด

\screenshot{admin-knowledge}{หน้าจัดการฐานความรู้แชตบอต}

\subsubsection{เมนูรายงานข้อมูลและคำขอกิจกรรม}

เมนู ``รายงานข้อมูล'' แสดงคิวรายงานข้อมูลสถานที่และกิจกรรมที่ผู้ใช้แจ้งว่าไม่ถูกต้อง ผู้ดูแลระบบตรวจสอบและดำเนินการ ดังภาพที่ \ref{fig:admin-reports} ส่วนเมนู ``คำขอกิจกรรม'' แสดงกิจกรรมที่สมาชิกส่งเข้ามา ผู้ดูแลระบบเลือก ``อนุมัติ'' เพื่อแสดงกิจกรรมบนหน้าเว็บ หรือ ``ไม่อนุมัติ'' พร้อมระบุเหตุผลที่ผู้ขอจะเห็น ดังภาพที่ \ref{fig:admin-event-requests}

\screenshot{admin-reports}{หน้ารายงานข้อมูล}

\screenshot{admin-event-requests}{หน้าคำขอกิจกรรม}

\subsubsection{เมนู QR Code และของรางวัล}

เมนู ``QR Code'' ใช้สร้าง QR Code ผูกกับสถานที่และกำหนดจำนวนพอยท์ พร้อมดูตัวอย่างและพิมพ์ป้าย ดังภาพที่ \ref{fig:admin-qr} ส่วนเมนู ``ของรางวัล'' ใช้เพิ่มและแก้ไขรายการของรางวัล ราคาพอยท์ และจำนวนคงเหลือ ดังภาพที่ \ref{fig:admin-rewards}

\screenshot{admin-qr}{หน้าจัดการ QR Code}

\screenshot{admin-rewards}{หน้าจัดการของรางวัล}

\subsubsection{เมนูแลกที่เคาน์เตอร์}

เมื่อสมาชิกมาแลกของรางวัล เจ้าหน้าที่เปิดเมนู ``แลกที่เคาน์เตอร์'' ดังภาพที่ \ref{fig:admin-redeem} แล้วดำเนินการ 2 ขั้นตอน คือ เลือกผู้ใช้ และเลือกของรางวัล จากนั้นกดยืนยันเมื่อพร้อมมอบของ ระบบจะหักพอยท์ตามราคาของรางวัล หากบันทึกผิดสามารถยกเลิกรายการจากประวัติได้ ซึ่งระบบจะคืนพอยท์ให้ผู้ใช้และคืนของรางวัล 1 ชิ้นกลับเข้าจำนวนคงเหลือ

\screenshot{admin-redeem}{หน้าแลกของรางวัลที่เคาน์เตอร์}

\subsubsection{เมนูผู้ใช้ ทริป และความพึงพอใจ}

เมนู ``ผู้ใช้'' ใช้ดูรายละเอียดและจัดการบัญชีผู้ใช้ ดังภาพที่ \ref{fig:admin-users} เมนู ``ทริป'' ใช้ค้นหาและดูแผนทริปที่ผู้ใช้สร้างไว้ ดังภาพที่ \ref{fig:admin-trips} และเมนู ``ความพึงพอใจ'' แสดงผลแบบสอบถามความพึงพอใจ ได้แก่ คะแนนเฉลี่ย จำนวนผู้ตอบ ข้อเสนอแนะ และคำตอบรายคน ดังภาพที่ \ref{fig:admin-survey}

\screenshot{admin-users}{หน้าจัดการผู้ใช้}

\screenshot{admin-trips}{หน้าดูทริปของผู้ใช้}

\screenshot{admin-survey}{หน้าผลความพึงพอใจ}
